\section{Empirical Strategy}
The empirical analysis is organized into two distinct parts.

First, using judicial decisions and administrative requests that obligate the government to provide free prescription and nonprescription drugs to individuals and restrict the government’s ability to plan and bargain with suppliers as a policy experiment, we compare ordinary and urgent purchases to estimate differences in the tender outcomes.

Finally, we estimate the “under the gun” effect, which, controlling for all other variables, consists of efficiency loss in public purchases due to the threat of punishing public officials in litigated tenders. In this case, we compare litigated and administrative purchases, two distinct types of urgent tenders, to identify the effect.

\subsection{The Enforcement Costs of Health Litigation}
As mentioned before, when there is an injunction or an administrative request that forces the government to make an urgent acquisition, the SES/SP has significantly less time to plan it and a lower degree of discretion in setting key procurement parameters than in an ordinary purchase. This first estimation aims to illustrate the impact of these exogenous and anomalous requests on the critical parameters of tenders.

First, it is essential to note that both judicial and administrative requests constitute shocks that are not correlated with any unobserved factors in the purchase process. Although a judicial or administrative order’s success depends on individuals’ characteristics, the order purchase process does not.

The principle of impersonality in public administration makes the planning and execution of the purchase utterly independent of the requesting individual’s characteristics. Thus, the purchase results do not depend on who placed the order or who will benefit from it. The tender outcomes depend on the purchase characteristics, such as the items to be purchased, planning, and market conditions.

Who makes the purchase, i.e., public officials of the PBUs, judicial orders, or administrative requests, functions as an exogenous restriction on the way they make purchases. Primarily, in this particular case, the shocks separate the purchases into two types according to planning conditions and required delivery time: ordinary and urgent purchases. Since planning conditions and delivery time are very similar between purchases based on judicial orders and administrative requests, they are both classified as urgent.

Thus, we identify the effects of these exogenous shocks on the tender results, comparing ordinary and urgent purchase types. Differences in reference prices between urgent and ordinary purchases of purchase order $i$, with a good $g$ and in time $t$, for instance, are estimated in the following specification for the log of reference price:

\[
\ln(\text{RefPrice}_{i,g,t}) = \alpha + \beta \cdot \text{Urgent}_{i,g,t} + \gamma_{i} + \lambda_{t} + \mathbf{X}_{i,g,t} + \varepsilon_{i,g,t},
\]

where $\gamma_{i}$ and $\lambda_{t}$ are item fixed effects and time trend dummies, respectively, and $\text{Urgent}$ is a dummy with value 1 if it is a litigated or administrative purchase and 0 if it is an ordinary purchase.

We use data of all public bids related to SUS-list medicines from January 2009 to December 2019. The data include only items with at least one urgent purchase and at least one ordinary tender. The results are shown in Table~4.

\begin{table}[ht]
  \centering
  \caption{Reference Prices: Urgent vs. Ordinary Purchases}
  \begin{tabular}{lcccc}
    \hline
           & (1)        & (2)        & (3)        & (4)        \\
           & OLS       & FE         & FE         & FE         \\
    \hline
    urgent     & 0.4988*** & 0.5243*** & 0.4967*** & 0.4725*** \\
               & (0.0379)  & (0.0383)  & (0.0377)  & (0.0373)  \\
    type\_mgmt & 0.6535*** &           & 0.6488*** & 0.553***  \\
               & (0.0516)  &           & (0.0510)  & (0.0521)  \\
    sealed-bid &           &           &           & -0.1708*** \\
               &           &           &           & (0.0184)  \\
    \_cons     & 0.33***   & 0.2554*** & -0.3851*** & -0.2307*** \\
               & (0.0607)  & (0.0699)  & (0.0807)  & (0.0809)  \\
    Observations & 59708   & 59708     & 59708     & 59708     \\
    R-squared  & 0.7767    & 0.0702    & 0.0781    & 0.0816    \\
    Item dummies & YES     & YES       & YES       & YES       \\
    Year dummies & NO      & YES       & YES       & YES       \\
    PBU dummies & YES      & NO        & YES       & YES       \\
    \hline
    \multicolumn{5}{l}{\footnotesize Standard errors are in parentheses. *** $p<0.01$, ** $p<0.05$, * $p<0.1$}
  \end{tabular}
\end{table}

It is possible to observe that reference prices are consistently higher in urgent purchases, considering all specifications: items are, on average, from 60.40\% to 68.93\% more expensive than in ordinary purchases.

The positive difference in reference prices captures the effects of worse conditions for planning purchases: smaller quantities, shorter delivery times, and the potential risk of punishment. As shown in Table~5, the quantities chosen are indeed lower for urgent purchases.

\begin{table}[ht]
  \centering
  \caption{Quantities: Urgent vs. Ordinary Purchases}
  \begin{tabular}{lcccc}
    \hline
           & (1)        & (2)        & (3)        & (4)        \\
           & OLS       & FE         & FE         & FE         \\
    \hline
    urgent     & -0.8128*** & -0.8811*** & -0.8272*** & -0.9402*** \\
               & (0.0625)   & (0.0623)   & (0.0626)   & (0.0635)   \\
    type\_mgmt & -1.2808*** &           & -1.2642*** & -1.7127*** \\
               & (0.1076)   &           & (0.1070)   & (0.1082)   \\
    sealed-bid &            &           &           & -0.800***  \\
               &            &           &           & (0.0428)   \\
    \_cons     & 7.412***   & 6.1144*** & 7.3626*** & 8.0859***  \\
               & (0.1117)   & (0.1003)   & (0.1458)   & (0.1492)   \\
    Observations & 59708    & 59708     & 59708     & 59708     \\
    R-squared  & 0.4141     & 0.1606    & 0.1702    & 0.1950    \\
    Item dummies & YES      & YES       & YES       & YES       \\
    Year dummies & NO       & YES       & YES       & YES       \\
    PBU dummies & YES       & NO        & YES       & YES       \\
    \hline
    \multicolumn{5}{l}{\footnotesize Standard errors are in parentheses. *** $p<0.01$, ** $p<0.05$, * $p<0.1$}
  \end{tabular}
\end{table}

Usually, governments seek to buy goods and services from the private sector on a large scale to obtain higher discounts on negotiated prices. This “bulk procurement” effect might be maximized if PBUs could adequately plan the acquisition process of goods and services.

The results suggest the following general mechanism: court orders or administrative orders create worse conditions for planning a given public tender. This shock impacts the capacity to define the amount to be purchased and the maximum prices PBUs are willing to pay.

Given that budgetary resources are scarce, especially for urgent orders, PBUs choose smaller quantities and pay higher prices to comply with court orders in the proper time. Thus, PBUs lose bargaining power and the possibility of substantial bulk discounts.

The effective fulfillment of court orders or administrative requests occurs in the external phase, consisting of the negotiation itself. Compliance with these external requests directly impacts the total amount of public spending. Unplanned, extrabudgetary resources are used to meet these external requests.

In urgent purchases, the negotiated quantities are, on average, nearly 53\% smaller than in ordinary purchases. Both reference prices and negotiated prices may be affected: the lower the quantities purchased, the higher the prices are.

We use a specification similar to that presented above, as a baseline to model differences in outcomes $y$ between urgent and ordinary purchases of purchase order $i$ with good $g$ and in time $t$:

\[
y_{i,g,t} = \alpha + \beta \cdot \text{Urgent}_{i,g,t} + \gamma_{i} + \lambda_{t} + \mathbf{X}_{i,g,t} + \varepsilon_{i,g,t},
\]

where $\mathbf{X}_{i,g,t}$ is a vector of controls, including the purchased quantities defined in the internal phase. This is a way to capture possible bulk discounts. The estimations for negotiated prices are presented in Table~6.

\begin{table}[ht]
  \centering
  \caption{Negotiated Prices: Urgent vs. Ordinary Purchases}
  \begin{tabular}{lcccc}
    \hline
           & (1)        & (2)        & (3)        & (4)        \\
           & OLS       & FE         & FE         & FE         \\
    \hline
    urgent     & 0.3672*** & 0.3526*** & 0.3568*** & 0.2680*** \\
               & (0.0414)  & (0.0429)  & (0.0417)  & (0.0432)  \\
    lquantity  & -0.3081*** & -0.3011*** & -0.3025*** & -0.3301*** \\
               & (0.0360)   & (0.0334)   & (0.0340)   & (0.0340)   \\
    type\_mgmt & -0.1814**  &           & -0.1422*  & -0.4735*** \\
               & (0.0801)  &           & (0.0742)  & (0.0863)   \\
    sealed-bid &           &           &           & -0.5403*** \\
               &           &           &           & (0.0344)   \\
    \_cons     & 2.4156*** & 1.2356*** & 1.3862*** & 2.0863*** \\
               & (0.2934)  & (0.2060)  & (0.2665)  & (0.2892)  \\
    Observations & 38440   & 38440     & 38440     & 38440     \\
    R-squared  & 0.8790    & 0.3393    & 0.3396    & 0.3798    \\
    Item dummies & YES     & YES       & YES       & YES       \\
    Year dummies & NO      & YES       & YES       & YES       \\
    PBU dummies & YES      & NO        & YES       & YES       \\
    \hline
    \multicolumn{5}{l}{\footnotesize Standard errors are in parentheses. *** $p<0.01$, ** $p<0.05$, * $p<0.1$}
  \end{tabular}
\end{table}

As shown, the government buys the same product under different planning situations: ordinary and urgent conditions. Negotiated prices are consistently higher for urgent purchases. On average, items purchased in adverse conditions are from 30.73\% to 44.37\% more expensive than those purchased in ordinary tenders.

Higher prices for urgent purchases suggest that adverse trading conditions strongly affect government bargaining power. On the other hand, tight deadlines and small quantities can alienate firms potentially interested in selling to the government. Table~7 presents estimations for the number of participant firms in urgent vs. ordinary purchases.

\begin{table}[ht]
  \centering
  \caption{Participant Firms: Urgent vs. Ordinary Purchases}
  \begin{tabular}{lcccc}
    \hline
           & (1)        & (2)        & (3)        & (4)        \\
           & OLS       & FE         & FE         & FE         \\
    \hline
    urgent     & -0.3887*** & -0.3831*** & -0.3811*** & -0.3373*** \\
               & (0.0190)   & (0.0191)   & (0.0191)   & (0.0200)   \\
    lquantity  & 0.1480***  & 0.1475***  & 0.1468***  & 0.1604***  \\
               & (0.0068)   & (0.0068)   & (0.0068)   & (0.0079)   \\
    type\_mgmt & -0.0463*   &           & -0.0672**  & 0.0960***  \\
               & (0.0267)   &           & (0.0281)   & (0.0293)   \\
    sealed-bid &           &           &           & 0.2661***  \\
               &           &           &           & (0.0184)   \\
    \_cons     & -0.3167*** & -0.0293    & 0.0419     & -0.3030*** \\
               & (0.0632)   & (0.0563)   & (0.0701)   & (0.0841)   \\
    Observations & 38430   & 38430     & 38430     & 38430     \\
    R-squared  & 0.4691    & 0.2645    & 0.2647    & 0.2886    \\
    Item dummies & YES     & YES       & YES       & YES       \\
    Year dummies & NO      & YES       & YES       & YES       \\
    PBU dummies & YES      & NO        & YES       & YES       \\
    \hline
    \multicolumn{5}{l}{\footnotesize Standard errors are in parentheses. *** $p<0.01$, ** $p<0.05$, * $p<0.1$}
  \end{tabular}
\end{table}

Although the number of participants for urgent purchases is approximately 30\% lower, the number of bids is even smaller. On average, the number of valid bids falls from 39.40\% to 45.93\% in urgent purchases compared to ordinary ones.

This firm behavior may reflect the lack of incentives to be more aggressive in the context of urgent purchases. Since there are fewer participating companies and less bargaining power for PBUs, suppliers make less effort to lower prices.

Moreover, adverse conditions for purchasing medicines can generate another problem related to firm screening processes. When conditions are precarious, the bidding process might fail. There may be a lack of interest among suppliers, or the PBUs cannot obtain a reasonable price. Urgent purchases tend to be significantly more likely to fail than ordinary purchases, as shown in Table~9.

\begin{table}[ht]
  \centering
  \caption{Successful Tenders: Urgent vs. Ordinary Purchases}
  \begin{tabular}{lcccc}
    \hline
           & (1)        & (2)        & (3)        & (4)        \\
           & LOGIT     & LOGIT      & LOGIT      & LOGIT      \\
    \hline
    urgent     & -0.4871*** & -0.4871*** & -0.6667*** & -0.5471*** \\
               & (0.0252)   & (0.0252)   & (0.0290)   & (0.0301)   \\
    lquantity  & 0.2378***  & 0.2378***  & 0.2508***  & 0.2803***  \\
               & (0.0060)   & (0.0060)   & (0.0063)   & (0.0065)   \\
    type\_mgmt &            &            & 0.1628**   & 0.7560***  \\
               &            &            & (0.0695)   & (0.0730)   \\
    sealed-bid &            &            &           & 0.9603***  \\
               &            &            &           & (0.0305)   \\
    \_cons     & -1.8201*** & -1.8201*** & -2.1149*** & -2.9174*** \\
               & (0.3403)   & (0.3403)   & (0.3536)   & (0.3615)   \\
    Observations & 59672   & 59672     & 59672     & 59672     \\
    Pseudo $R^2$ & 0.1238  & 0.1238    & 0.1328    & 0.1572    \\
    Item dummies & YES     & YES       & YES       & YES       \\
    Year dummies & NO      & NO        & NO        & YES       \\
    PBU dummies & NO       & NO        & YES       & YES       \\
    \hline
    \multicolumn{5}{l}{\footnotesize Standard errors are in parentheses. *** $p<0.01$, ** $p<0.05$, * $p<0.1$}
  \end{tabular}
\end{table}

Urgent purchases are from 38.56\% to 48.66\% less likely to succeed than ordinary purchases. Failure to bid has relevant implications in terms of budgetary costs. First, when a court-ordered purchase is not made, it generates relevant punishment costs, such as fines and blocking of budgetary resources for the PBU. In addition, as mentioned in section 2.3, planning and executing a purchase have high costs. Therefore, resources are wasted in case of failure.

\subsection{The “Under the Gun” Effect}
As already mentioned, urgent purchases consist of those arising from court orders and administrative orders. Both litigated and administrative purchases are made under challenging conditions in terms of planning and execution. However, there is a single significant difference between them: a litigated tender likely results in a punishment for public officials if they fail to complete the purchase.

Making public the information that a tender is of a litigated type might have an additional effect on the results of bidding processes. Since all participants learn that the government is under even higher pressure to purchase, firms have additional advantages over PBUs in the bargaining process.

This section’s main objective is to estimate this “additional effect”; we call it the “under the gun” effect. We use all public bid data, including SUS-list and non-SUS-list medicines, from January 2009 to December 2019.

We restrict the analysis to items with at least one litigated purchase and at least one administrative tender; data on ordinary purchases are excluded from those estimations. The idea is to compare litigated and administrative purchases exclusively. Using the same identification strategy presented in section 4.1, we adopt the following specification as the baseline equation:

\[
\ln(\text{Price}_{i,g,t}) = \alpha + \beta \cdot \text{Admin}_{i,g,t} + \gamma_{i} + \lambda_{t} + \mathbf{X}_{i,g,t} + \varepsilon_{i,g,t},
\]

where $\text{Admin}$ has a value of 1 if it is an administrative purchase and 0 if it is a litigated purchase. Table~10 presents the estimation of the results for negotiated prices. Negotiated prices are lower for administrative purchases. On average, administrative tenders are 8.11\% to 9.06\% less expensive than litigated tenders.

\begin{table}[ht]
  \centering
  \caption{Negotiating Prices: The “Under the Gun” Effect, Litigated vs. Administrative Tenders}
  \begin{tabular}{lcccc}
    \hline
                 & (1)        & (2)        & (3)        & (4)        \\
                 & OLS       & FE         & FE         & FE         \\
    \hline
    administrative   & -0.095*** & -0.0859*** & -0.0859*** & -0.0846*** \\
                     & (0.0238)  & (0.0243)   & (0.0243)   & (0.0242)   \\
    lquantity        & -0.4003*** & -0.3942*** & -0.3942*** & -0.3981*** \\
                     & (0.0233)   & (0.0236)   & (0.0236)   & (0.0234)   \\
    type\_mgmt       & -0.1564   &           & -0.1840    & -0.4322*   \\
                     & (0.1810)  &           & (0.1879)   & (0.2519)   \\
    sealed-bid       &           &           &           & -0.5173*** \\
                     &           &           &           & (0.0401)   \\
    \_cons           & 2.4521*** & 3.9031*** & 4.0871*** & 4.5326*** \\
                     & (0.2431)  & (0.1344)   & (0.2391)   & (0.3013)   \\
    Observations     & 51013     & 51013     & 51013     & 51013     \\
    R-squared        & 0.9293    & 0.3718    & 0.3718    & 0.3837    \\
    Item dummies     & YES       & YES       & YES       & YES       \\
    Year dummies     & NO        & YES       & YES       & YES       \\
    PBU dummies      & YES       & NO        & YES       & YES       \\
    \hline
    \multicolumn{5}{l}{\footnotesize Standard errors are in parentheses. *** $p<0.01$, ** $p<0.05$, * $p<0.1$}
  \end{tabular}
\end{table}

In other words, litigated purchases are between 8.83\% and 9.97\% more expensive than administrative purchases. This difference is the “under the gun” effect. With similar planning and execution conditions between administrative and litigated purchases, the estimated price difference can be attributed exclusively to the possible punishment of PBUs in case of failure to purchase.
